\documentclass[a4paper,10pt,fleqn]{article}
\usepackage[margin=1in]{geometry}
\usepackage{amssymb}
\usepackage{adjustbox}
\usepackage{natbib}
\usepackage[colorlinks=true,linkcolor=blue,citecolor=blue,urlcolor=blue]{hyperref}
\usepackage{fuzz}
\usepackage{zed-maths}
\usepackage{zed-proof}
\newdimen\savedleftskip
\begin{document}

\section*{Schema Renaming}

\noindent Schema renaming lets you create a version of an existing schema with some component names changed.  The notation S[old/new] is per Z RM §3.11.

\bigskip

\begin{schema}{Counter}
count : \nat
\where
count \geq 0
\end{schema}

\noindent Single rename: produce a Counter with component ''count'' renamed to ''n''.

\bigskip

\begin{zed}
CounterN \defs Counter[count/n]
\end{zed}

\noindent Multiple renames in one step:

\bigskip

\begin{schema}{Point}
x : \nat \\
y : \nat
\end{schema}

\noindent Rename both components at once.

\bigskip

\begin{zed}
SwappedPoint \defs Point[x/y, y/x]
\end{zed}

\section*{Renaming Primed Schemas}

\noindent The decoration is tightest-binding per Z RM §3.11: S'[a/b] renames the primed schema S', not the result of renaming then priming.

\bigskip

\begin{zed}
Op2 \defs Counter[count'/count]
\end{zed}

\noindent In an operation schema context, this captures the after-state renaming: the primed component count' is renamed back to count in the view.

\bigskip

\section*{Decorated Component Names}

\noindent The rename pairs themselves may carry decoration.

\bigskip

\begin{zed}
CounterOp \defs Counter[count'/n]
\end{zed}

\end{document}